
\subsection*{State of the art and preliminary work}

\subsubsection*{Remote Sensing and Embeddings Representations}

\begin{figure}
  \includegraphics[width=\textwidth]{figures/embeddings}
  \caption{\textbf{Location Embedding differences capture policy discontinuities}. \textbf{A)} Embeddings originate from Natural Language Processing, where they semantically describe words (or sentences) in a metric vector space that allows for arithmetic operations: The difference $\Delta e$ between the words "King" and "Queen" capture the semantics of gender so that $e("man") + \Delta e = e("woman")$. \textbf{B)} Analogously, location embeddings semantically describe coordinates in a vector space that can capture socio-economic and environmental patterns at that location. \textbf{C)} Analyzing embeddings near borders, as on the shown china-russia example, factors out environmental factors and captures policy-induced differences like deforestation. A link to the proof of concept is available on this \href{https://code.earthengine.google.com/0550b5b8b19d8014888410b47f463231}{Google Earth Engine Script}}
   
\end{figure}

\noindent\textbf{From word embeddings to Earth embeddings.}
Embedding representations originated in Natural Language Processing (NLP) as dense vectors that capture semantic similarity and enable linear analogies in a metric space, turning discrete symbols into a representation suitable for statistical learning. This conceptual shift motivates \emph{Earth embeddings} as an analogous representation layer for geospatial coordinates: instead of describing words, location embeddings describe places by mapping longitude/latitude (and potentially radius/time) into a semantic vector space that supports similarity search, retrieval, and downstream modelling \citep{klemmer2025earth}. In this view, embeddings are not merely features, but a representation interface between heterogeneous data sources and statistical/econometric workflows.

\noindent\textbf{Two complementary paradigms: embedding databases vs.\ embedding models.}
Current Earth-embedding approaches fall into two complementary classes \citep{klemmer2025earth}. \emph{Embedding databases} store precomputed feature vectors for many locations derived from remote-sensing foundation models, i.e., explicit encoders applied to satellite imagery or time series. Early work used accessible global imagery pipelines for scalable ML \citep{rolf2021generalizable}, and recent systems increasingly rely on deep encoders (e.g., ResNets/Vision Transformers) to provide dense features at scale \citep{czerkawski2024majortom}. A prominent example is Google DeepMind’s AlphaEarth Foundations (AEF), which releases high-resolution, dense, pixel-wise embeddings that can be used as an analysis-ready representation layer for downstream tasks and spatial comparisons \citep{brown2025alphaearth}. This ``database'' view is practically attractive because it amortizes preprocessing and enables large-scale retrieval over massive archives \citep{klemmer2025earth}.

In contrast, \emph{embedding models} represent locations \emph{implicitly} via a learned \emph{location encoder} that maps coordinates to embeddings continuously, i.e., as an implicit neural representation (INR) indexed by coordinates \citep{mai2022review}. This model-based view is particularly natural for geospatial data, where coordinates provide the shared key that joins heterogeneous modalities \citep{klemmer2025earth}. Methodologically, it builds on two broad AI developments emphasized in the Emmy Noether proposal: (i) \emph{multimodal contrastive learning}, popularized by CLIP-style objectives \citep{radford2021clip}, and (ii) \emph{implicit neural representations} that store signals in neural network weights and interpolate smoothly between observations \citep{sitzmann2020siren,mildenhall2021nerf}. For geospatial embeddings, self-supervised \emph{geolocalization} provides the key training signal: the model learns an embedding that is predictive of where an observation came from \citep{klemmer2025earth}.

\noindent\textbf{Low-resolution global embedding models (SatCLIP) and multimodal foundation models.}
A representative example of an \emph{embedding model} is SatCLIP, which learns global, general-purpose location embeddings directly from satellite imagery using a CLIP-inspired objective, yielding a continuous location encoder that can be queried at arbitrary coordinates \citep{klemmer2025satclip,klemmer2025earth}. Similarly, GeoCLIP learns location embeddings via geolocalization using ground-level imagery \citep{vivanco2023geoclip}. More broadly, geospatial foundation models increasingly provide embeddings across sensors and time (e.g., optical/radar, multitemporal stacks), with recent systems such as Tessera emphasizing scalable multimodal, multiscale remote-sensing representations \citep{feng2025tessera,klemmer2025earth}. Together, these developments establish Earth embeddings as an emerging representation layer for geospatial sciences—much like word embeddings transformed NLP—while highlighting open methodological gaps in how embeddings are constructed, validated, and made useful for downstream scientific inference \citep{klemmer2025earth}. % (This mirrors the database-vs-model distinction and INR/multimodal framing in the Emmy Noether proposal.) 

\subsubsection*{Border Analysis policy-relevant Case Studies}

\noindent\textbf{Border designs: using administrative discontinuities for credible causal inference.}
Borders provide a powerful quasi-experimental setting because many policies change discretely at administrative boundaries while biophysical conditions often vary smoothly in space. This motivates regression discontinuity designs (RDD), which identify local treatment effects by comparing outcomes on either side of a cutoff (here: distance to the border), under continuity assumptions for potential outcomes and covariates \citep{imbenslemieux2008}. In the remote-sensing context, the outcome can be a proxy indicator (e.g., NDVI) or an embedding-derived index (e.g., distance in embedding space), evaluated for locations within a bandwidth around borders and modelled as a function of signed distance to the boundary.

\noindent\textbf{RDD with Earth observation: from proxies to embedding outcomes.}
W\"upper et al.\ demonstrate the approach for forest-policy effectiveness by sampling remote-sensing indicators near many international borders and estimating the discontinuity gap as evidence of differential policy impacts \citep{Wuepper2024GlobalForestBorders}. This setting is naturally suited to embedding representations: if embeddings capture land management and land-cover dynamics, then \emph{discontinuities in embedding space} can serve as a rich outcome representation that potentially increases sensitivity to subtle changes beyond single-index proxies.

\noindent\textbf{Differences-in-discontinuities and border DID.}
Many policy questions require distinguishing persistent cross-border differences from \emph{policy-induced changes over time}. A practical extension is to combine border RDD with time variation, i.e., a \emph{difference-in-discontinuities} design that contrasts the border discontinuity before vs.\ after a policy change (or across policy regimes), thereby differencing out time-invariant border-specific confounding \citep{grembi2016fiscal}. Complementarily, difference-in-differences (DID) compares changes over time between treated and control units under a (conditional) parallel-trends assumption; modern DID estimators explicitly address multiple time periods and staggered adoption, which is common in environmental policy \citep{callaway2021did}. In this proposal’s case studies, these designs translate into: (i) spatial local comparisons around borders (RDD), and (ii) temporal contrasts of border gaps (differences-in-discontinuities) or panel-style comparisons where treatment timing varies across regions (DID), with Earth embeddings providing a unified measurement layer across locations and years.

